%  Short title: the full one overflows the running header.
\chapter[Mass Storage: SD Cards, SPI Flash, ext4]%
        {Mass Storage: SD Cards, SPI Flash, and the ext4 Filesystem}
\label{ch:storage}

The final layer is mass storage: block drivers for SD cards and for a raw SPI NOR
flash chip, and a complete pure-Ada ext2/3/4 filesystem (a reimplementation of
lwext4) that sits on top of them: mounting real \texttt{mkfs.ext4} cards,
\emph{and} wear-leveling and formatting a bare flash chip entirely on-device. The
smallest non-volatile stores --- serial EEPROM and its faster, wear-free cousin
FRAM, each driven by a generic over a whole catalogue of parts --- get their own
chapters next (Chapters~\ref{ch:eeprom} and~\ref{ch:fram}).

\section{SD\_SPI: an SD card over the SPI master}\index{SD card}\index{SD\_SPI}\index{SPI SD mode}

\texttt{ESP32S3.SD\_SPI} talks to an SD/SDHC card in SPI mode, layered on the
task-safe SPI driver. The chip-select is an ordinary GPIO so it can be held across
a whole command. Blocks are 512-byte logical sectors (LBA addressing); SDSC and
SDHC are detected and handled internally.

\begin{lstlisting}
type Block is array (0 .. 511) of Interfaces.Unsigned_8;
type Status is (OK, No_Card, Unusable, Init_Timeout, Read_Error, Write_Error);
procedure Setup (C : out Card; Host : ESP32S3.SPI.SPI_Host;
                 Sclk, Mosi, Miso, Cs : Pin_Id;
                 Init_Clock_Hz : Positive := 400_000;
                 Data_Clock_Hz : Positive := 8_000_000);
procedure Initialize (C : in out Card; Result : out Status);
procedure Read_Block  (C : in out Card; LBA : Block_Address;
                       Data : out Block; Result : out Status);
procedure Write_Block (C : in out Card; LBA : Block_Address;
                       Data : Block; Result : out Status);
\end{lstlisting}

\lstexample{Bring up a card and read sector 0.}
\begin{lstlisting}
with ESP32S3.SD_SPI; use ESP32S3.SD_SPI;
with ESP32S3.SPI;
procedure SD_Init is
   C   : Card;
   St  : Status;
   Sec : Block;
begin
   Setup (C, ESP32S3.SPI.SPI2, Sclk => 12, Mosi => 11, Miso => 13, Cs => 10);
   Initialize (C, St);
   if St = OK then
      Read_Block (C, LBA => 0, Data => Sec, Result => St);   --  the MBR sector
   end if;
end SD_Init;
\end{lstlisting}

\lstexample{A non-destructive read / write-back / re-read round-trip.}
\begin{lstlisting}
with ESP32S3.SD_SPI; use ESP32S3.SD_SPI;
with ESP32S3.SPI;
procedure SD_Roundtrip is
   C          : Card;
   St         : Status;
   Orig, Back : Block;
begin
   Setup (C, ESP32S3.SPI.SPI2, Sclk => 12, Mosi => 11, Miso => 13, Cs => 10);
   Initialize (C, St);
   Read_Block  (C, 16#2000#, Orig, St);     --  read a scratch sector
   Write_Block (C, 16#2000#, Orig, St);     --  write the same bytes back
   Read_Block  (C, 16#2000#, Back, St);     --  re-read; Back = Orig, no data lost
end SD_Roundtrip;
\end{lstlisting}

\section{SDMMC: the native SD/MMC host}\index{SDMMC}\index{SDHOST controller}

\texttt{ESP32S3.SDMMC} drives the dedicated SDHOST controller (clock + command +
1- or 4-bit data) in PIO mode, which is faster than SPI mode, and with no DMA or
finalization it works under every profile. The API mirrors \texttt{SD\_SPI}.

\begin{lstlisting}
procedure Setup (C : out Card; On : Slot; Clk, Cmd, D0 : Pin_Id;
                 D1, D2, D3 : Optional_Pin := No_Pin;
                 Width : Bus_Width := Width_1; ...);
procedure Initialize (C : in out Card; Result : out Status);
procedure Read_Block  (C : in out Card; LBA : Block_Address; Data : out Block;
                       Result : out Status);
\end{lstlisting}

\lstexample{4-bit bus, read a sector.}
\begin{lstlisting}
with ESP32S3.SDMMC; use ESP32S3.SDMMC;
procedure SDMMC_Read is
   C   : Card;
   St  : Status;
   Sec : Block;
begin
   Setup (C, On => Slot1, Clk => 14, Cmd => 15, D0 => 2, D1 => 4, D2 => 12,
          D3 => 13, Width => Width_4);
   Initialize (C, St);
   if St = OK then
      Read_Block (C, 0, Sec, St);
   end if;
end SDMMC_Read;
\end{lstlisting}

\section{The ext4 filesystem}\index{ext4 filesystem}\index{lwext4}\index{JBD2 journal}

\texttt{ESP32S3.Ext4} is a pure-Ada ext2/3/4 implementation with a JBD2 journal
(Chapter overview in \texttt{libs/esp32s3\_hal/src/ext4/README.md}). It reads
default \texttt{mkfs.ext4} cards (extents, 64-bit, HTree directories,
\texttt{metadata\_csum} verification), writes ext2/3/non-csum filesystems, and
recovers a dirty journal on mount. It needs the \texttt{embedded} or \texttt{full}
profile.

\subsection{Connecting a block device}\index{block device}

The filesystem talks to a small block-device abstraction (\texttt{ESP32S3.%
Block\_Dev.Device}). An adapter turns an initialised SD card into one with a
single call. The \texttt{Make} below appears in each complete example that
follows:

\begin{lstlisting}
with ESP32S3.Block_Dev;
with ESP32S3.Block_Dev.SD_SPI_Source;
--  ... given an initialised "Card : aliased ESP32S3.SD_SPI.Card" ...
   Dev : constant ESP32S3.Block_Dev.Device :=
           ESP32S3.Block_Dev.SD_SPI_Source.Make (Card'Access);
\end{lstlisting}

\subsection{Mounting and reading}\index{mount}\index{inode}\index{journal replay}

A \texttt{Mount} is a controlled handle: it flushes and closes on scope exit, and
\texttt{Open} automatically replays a dirty journal on a writable mount. The
complete example brings up the card, wraps it, mounts read-only, then reads a
file:

\begin{lstlisting}
with ESP32S3.SPI;
with ESP32S3.SD_SPI;
with ESP32S3.Block_Dev;
with ESP32S3.Block_Dev.SD_SPI_Source;
with ESP32S3.Ext4;
with ESP32S3.Ext4.FS;
with ESP32S3.Ext4.Inode;
procedure EXT4_Read is
   use type ESP32S3.SD_SPI.Status;
   Card : aliased ESP32S3.SD_SPI.Card;
   St   : ESP32S3.SD_SPI.Status;
begin
   ESP32S3.SD_SPI.Setup (Card, ESP32S3.SPI.SPI2,
                         Sclk => 12, Mosi => 11, Miso => 13, Cs => 10);
   ESP32S3.SD_SPI.Initialize (Card, St);
   if St /= ESP32S3.SD_SPI.OK then
      return;
   end if;
   declare
      Dev  : constant ESP32S3.Block_Dev.Device :=
               ESP32S3.Block_Dev.SD_SPI_Source.Make (Card'Access);
      M    : ESP32S3.Ext4.FS.Mount;
      I    : ESP32S3.Ext4.Inode.Info;
      Buf  : ESP32S3.Ext4.Byte_Array (0 .. 255);
      Last : Natural;
   begin
      M.Open (Dev, Read_Only => True);
      M.Stat (M.Lookup ("/etc/hello.txt"), I);  --  resolve a path -> its inode
      M.Read_File (I, Offset => 0, Into => Buf, Last => Last);
      --  Buf (0 .. Last - 1) holds the first Last bytes of the file
   end;   --  M finalizes: flush + close
end EXT4_Read;
\end{lstlisting}

\textbf{Listing a directory.} \texttt{Iterate} calls a visitor for each entry.
The card bring-up is identical to the example above; the body is:

\begin{lstlisting}
   declare
      Dev  : constant ESP32S3.Block_Dev.Device :=
               ESP32S3.Block_Dev.SD_SPI_Source.Make (Card'Access);
      M    : ESP32S3.Ext4.FS.Mount;
      Root : ESP32S3.Ext4.Inode.Info;
      procedure Show (Name : String;
                      Ino  : ESP32S3.Ext4.Inode_Number;
                      File_Type : ESP32S3.Ext4.U8) is
      begin
         Put_Line ("  " & Name);
      end Show;
   begin
      M.Open (Dev, Read_Only => True);
      M.Stat (M.Lookup ("/"), Root);
      M.Iterate (Root, Show'Access);
   end;
\end{lstlisting}

\subsection{Writing}

On a non-\texttt{metadata\_csum} filesystem (ext2/3, or
\texttt{mke2fs -O \textasciicircum metadata\_csum}) the filesystem can be modified.
The full POSIX-like set is available: \texttt{Create\_File}, \texttt{Write\_File},
\texttt{Truncate}, \texttt{Mkdir}, \texttt{Rmdir}, \texttt{Unlink},
\texttt{Rename} and hard \texttt{Link}.

\begin{lstlisting}
   declare
      Dev : constant ESP32S3.Block_Dev.Device :=
              ESP32S3.Block_Dev.SD_SPI_Source.Make (Card'Access);
      M   : ESP32S3.Ext4.FS.Mount;
      N   : ESP32S3.Ext4.Inode_Number;
   begin
      M.Open (Dev, Read_Only => False);
      N := M.Create_File ("/", "log.txt");
      M.Write_File (N, (Character'Pos ('h'), Character'Pos ('i'), 10));
      M.Mkdir ("/", "data");
      M.Close;          --  flush + close (or let M finalize)
   end;
\end{lstlisting}

\subsection{Streaming writes and file-size limits}\index{streaming write}\index{Append}\index{block map}\index{single-indirect block}\index{double-indirect block}

\texttt{Write\_File} sets a file's whole contents from one in-memory buffer.
That is fine for a small file, but you cannot hand an embedded part a gigabyte-sized
array. \texttt{Append} grows a file at its end one chunk at a time, so a large
file is built from small writes with no buffer to hold it all:

\begin{lstlisting}
   N := M.Create_File ("/", "big.log");
   for I in 1 .. 1000 loop          --  e.g. 1000 records, streamed
      M.Append (N, Chunk);          --  grows the file; no whole-file buffer
   end loop;
   M.Commit;
\end{lstlisting}

What bounds a single file is the \emph{block map}. The writer allocates classic
indirect blocks (12 direct, then single-indirect, then double-indirect), so
on a 4\,KiB-block volume one file reaches about 4\,MiB through the single-indirect
tier and about 4\,GiB through the double-indirect tier; the matching free and
truncate paths walk the same tiers, so a file that can be written can always be
deleted and shrunk. The \emph{read} path goes further: it follows the full
classic map (through triple-indirect) and ext4 extent trees, so a
multi-gigabyte file written by a host is readable on-device, bounded only by the
volume's free space.

\subsection{Directory capacity}\index{directory capacity}\index{dirent}\index{directory growth}

A directory is itself a file of variable-length entries (each 8 bytes plus the
name, rounded to 4), and creating a name must find room for one. When every
existing directory block is packed, \texttt{Add\_Entry} grows the directory:
it allocates a block, writes a single spanning unused entry (the shape
\texttt{mkfs} gives a fresh directory block), wires it into the inode's next
\emph{direct} map slot, and persists the grown inode before placing the name.
Twelve direct slots at 4\,KiB give 48\,KiB of entries -- on the order of
1{,}500--2{,}900 names depending on their length. Growth past the direct slots
raises \texttt{No\_Space}, and a directory a host created with an extent map
raises \texttt{Unsupported\_Feature} rather than being modified half-correctly.

That ceiling is not hypothetical bookkeeping. Before growth existed, a
directory was capped at its \emph{first} block -- roughly 140 mirror-length
names -- and one application built on this SDK hit it silently: it mirrors a
weather service's chart set at three files per product, filled the block
partway through the first full download, and every later create failed with
\texttt{No\_Space} from that day on. The products that happened to sort late in
its list simply never appeared on the card, which surfaced to the user as one
category of the UI ``only ever shows one image.'' A capacity limit in a
filesystem rarely announces itself where the limit lives; budget directory
sizes (files per directory, not just bytes per volume) when planning on-card
layouts, and prefer one subdirectory per product family once counts head into
the hundreds.

\subsection{Crash-safe transactions}\index{crash-safe transactions}\index{journaling}\index{commit barrier}

For atomicity, group an operation behind \texttt{Commit}: the dirty metadata and
the superblock are written to the journal (durably), the RECOVER flag is set
(the commit barrier), then the metadata is checkpointed to its final locations
and the flag cleared. A power loss after the barrier is recovered forward by the
next \texttt{Open} (or by any other ext4 implementation, including the Linux
kernel); a loss before it has no effect.

\begin{lstlisting}
   --  Payload : constant ESP32S3.Ext4.Byte_Array := (16#41#, 16#42#, 16#43#);
   N := M.Create_File ("/", "important.txt");
   M.Write_File (N, Payload);
   M.Commit;         --  one atomic, journaled, crash-safe transaction
\end{lstlisting}

A volume formatted \emph{without} a journal (\texttt{mke2fs -O \textasciicircum
has\_journal}) has no log to write through, so \texttt{Commit} instead flushes the
dirty metadata and superblock straight to their final locations. The filesystem
detects which kind it mounted and routes \texttt{Commit} automatically. A
no-journal volume trades the crash-safety barrier for smaller size and fewer
writes, the right choice on a small flash chip.\index{no-journal volume}

\subsection{Concurrency: one owner per mount}\index{block cache}\index{single-owner mount}

A \texttt{Mount} is a \emph{single-owner} handle, not a shared service. Every
operation takes it \texttt{in out} and mutates shared state, and the filesystem
holds no internal lock or protected object: drive one mount from one task, and let
its calls run to completion in sequence.

That rule is stricter than it first looks because even reads have side effects.
Underneath the mount sits a write-back LRU block cache, and \texttt{Read\_File}
fetches-on-miss and reorders it, so two tasks ``only reading'' the same mount
still race on the cache. Writes are stronger still: \texttt{Create\_File},
\texttt{Write\_File}, \texttt{Mkdir} and the rest accumulate a dirty-metadata
write-set that \texttt{Commit} flushes as \emph{one} journaled transaction (the
crash-safe commit above). Interleaving two tasks' writes scrambles that write-set
and breaks the crash-safety guarantee, which assumes a single writer building one
coherent transaction. The block drivers \emph{beneath} the filesystem
(\texttt{SD\_SPI}, \texttt{SDMMC}) are themselves task-safe, but that only
serialises the \emph{bus}; it does nothing for the filesystem state above it.

If more than one task must reach the filesystem, give the mount a single owner and
funnel every request through it: a dedicated task that owns the \texttt{Mount} and
performs the I/O, with other tasks handing it work through a \texttt{protected}
request queue (Ravenscar-compatible) or, on the \texttt{full} profile, a rendezvous
entry. The I/O itself must stay in the owning task, never inside a protected
action because the SD drivers can suspend and a protected operation must not block.

\begin{lstlisting}
--  One task owns the Mount and is the only code that touches it; other tasks
--  hand it work (a protected request queue under Ravenscar, or a rendezvous
--  entry on the full profile).  The blocking I/O stays in this task.
task body File_Server is
   M : ESP32S3.Ext4.FS.Mount;          --  the single owner
begin
   M.Open (Dev, Read_Only => False);
   loop
      --  take the next request, then service it on the one mount:
      --     M.Read_File (...)              --  a read request
      --     M.Write_File (...);  M.Commit; --  a write request, committed
      null;
   end loop;
end File_Server;
\end{lstlisting}

\subsection{Errors}

Filesystem operations raise: the standard \texttt{Ada.IO\_Exceptions} names
(\texttt{Name\_Error} for a missing path, \texttt{Use\_Error}, \texttt{Device\_%
Error}, \texttt{End\_Error}) plus filesystem-specific \texttt{Unsupported\_%
Feature}, \texttt{Bad\_Checksum}, \texttt{Corrupt}, \texttt{No\_Space},
\texttt{Not\_Empty} and \texttt{Read\_Only}. Wrap calls in a handler so a bad card
or a missing file never escapes a task:

\begin{lstlisting}
begin
   M.Stat (M.Lookup ("/maybe/missing"), I);
exception
   when ESP32S3.Ext4.Name_Error => Put_Line ("not found");
   when ESP32S3.Ext4.Bad_Checksum => Put_Line ("corrupt metadata");
end;
\end{lstlisting}

\subsection{Free-count integrity: the phantom-free tripwire}\index{phantom-free tripwire}\index{free-count drift}

A subtle class of corruption is the \emph{free-count drift}: the superblock and
block-group descriptors cache how many blocks and inodes are free, and if a
\texttt{Free} adjusts that cache for a bit that was \emph{already} clear (a
double-free, or a stale/incoherent bitmap read off a flaky SD card), the cached
count creeps above what the bitmap actually says. \texttt{e2fsck} then reports a
Pass-5 count mismatch on an otherwise-sound filesystem.

The fix makes \texttt{Free} \emph{idempotent} by construction. \texttt{Clear\_Bit}
now reports whether the bit was actually set, and the group/superblock counts are
bumped \emph{only} when it was, so a double or phantom free is a no-op rather
than a drift, and the count cannot diverge from the bitmap.

\begin{lstlisting}
--  Bitmap.Clear_Bit returns False if the bit was already clear:
if Bitmap.Clear_Bit (Group_BM, Index) then
   Dec_Used (Group);          --  bit really was set: adjust the counts
else
   Phantom_Frees := Phantom_Frees + 1;   --  already free: no drift, but flag it
end if;
\end{lstlisting}

Crucially the already-clear path is \emph{not} silent. That would mask a genuine
double-free bug. A \texttt{Phantom\_Free\_Count} tripwire counts the events (a plain
counter, not a \texttt{pragma Assert} because the target also builds with
\texttt{-gnata} and an assert would abort the filesystem operation). The host test
harness asserts the count is zero after every scenario, and the on-target
\texttt{esp32s3\_ext4\_write} example reports it as a battery check
(\texttt{[PASS] no phantom frees}; a flaky card flips it to \texttt{[FAIL]}).

\begin{lstlisting}
Bitmap.Reset_Phantom_Free_Count;            --  arm the tripwire
--  ... exercise allocate / free / truncate / unlink ...
Check ("no phantom frees (count consistent with bitmap)",
       Bitmap.Phantom_Free_Count = 0);
\end{lstlisting}

The original drift was ultimately traced to a \emph{worn} SD card: a brand-new card
runs the whole battery \texttt{e2fsck}-clean with zero phantom frees. The tripwire
is what turned an intermittent, card-dependent count mismatch into a single
yes/no signal, surfacing the fault instead of hiding it.

\subsection{How it is validated}

Because the filesystem is pure logic over the block interface, it is developed and
tested host-native (x86) against real \texttt{mke2fs} images, with the Linux
kernel and \texttt{e2fsck} as the oracle: every read is byte-compared to the
source file, every mutation is checked with \texttt{e2fsck -fn}, and journal
replay and commit are cross-checked against the kernel's own recovery. The same
code cross-compiles for the ESP32-S3 and runs over the SD drivers above. See
\texttt{tests/ext4\_host/} and \texttt{libs/esp32s3\_hal/src/ext4/REFERENCES.md}.

\subsection{Scope: deliberately smaller than Linux ext4}\index{ext4 scope}\index{HTree}\index{delayed allocation}\index{metadata\_csum}

It is worth being explicit that this filesystem does not aim for feature parity
with the Linux kernel's ext4, and why. The asymmetry is intentional: the
\emph{read} path is wide (extents, 64-bit, HTree-indexed directories,
\texttt{metadata\_csum} verification -- a default desktop card mounts and reads),
while the \emph{write} path is a narrow, fully-understood subset (classic
indirect maps, linear directories, non-csum volumes). Anything outside that
subset is \emph{refused} -- \texttt{Read\_Only}, \texttt{Unsupported\_Feature} --
rather than attempted, because the most dangerous state for an on-disk format
is ``implemented enough to write, not enough to write correctly'': that is
silent corruption a desktop mount discovers weeks later. Refusal fails loudly
at the call site.

The features left out earn nothing at this scale. HTree directory indexing
beats a linear scan somewhere in the thousands of entries; a directory here
tops out around two thousand entries in 48\,KiB that is hot in the block cache.
Delayed allocation and the multi-block allocator exist to optimise allocation
locality under multi-process write pressure that a single-owner embedded mount
does not have. Full \texttt{jbd2} parity would double metadata writes on a bus
measured in single-digit MB/s (and add SD wear) to strengthen guarantees the
simple journal-plus-\texttt{Commit} model already provides. And every one of
them wants cache and working state sized for a kernel page cache, not for a
32\,KiB block cache sharing PSRAM with an application's framebuffers.

What keeps the small subset honest is that it stays small enough to
\emph{verify}: the host harness runs whole scenario batteries with
\texttt{e2fsck} as the referee, as described above, and individual pieces are
within reach of SPARK proof. A port of Linux ext4's roughly sixty thousand lines
would be the largest, least-reviewable component in the SDK, with a test
matrix that grows combinatorially in feature interactions -- correctness that
upstream earned through two decades of fleet exposure no embedded project can
replicate. If a concrete interoperability need arises -- the strongest
candidate is \emph{writing} \texttt{metadata\_csum} volumes, so cards formatted
by a modern desktop \texttt{mkfs.ext4} need not be reformatted -- the right
move is to implement that one feature (checksum recompute on the metadata the
writer already touches), not to chase parity.

\section{Putting ext4 on raw SPI NOR flash}\index{SPI NOR flash}\index{wear leveling}\index{flash translation layer}

The same filesystem runs on a bare SPI NOR flash chip, no SD card in sight. Four
pure-Ada layers stack up, each a thin adapter over the one below:

\begin{center}
\texttt{ESP32S3.Ext4} $\rightarrow$ \texttt{Block\_Dev.WL} $\rightarrow$
\texttt{Block\_Dev.W25Q\_Source} $\rightarrow$ \texttt{ESP32S3.W25Q}
\end{center}

\noindent The flash driver speaks to the chip; \texttt{W25Q\_Source} presents it
as a 512-byte block device; \texttt{WL} is a wear-leveling flash translation
layer that spreads writes so no sector wears out first; and \texttt{Ext4} mounts
the result. The reference part is a Winbond \texttt{W25Q256FV} (32\,MB) that shares
the SPI2 bus with the W5500 Ethernet chip.

\subsection{The W25Q flash driver}\index{W25Q}\index{JEDEC ID}\index{4-byte addressing}\index{chip-select}

\texttt{ESP32S3.W25Q} reads the JEDEC id, then programs, erases and reads in
\emph{4-byte address mode} (a part larger than 16\,MB needs four address bytes).
Two design points stand out. The chip size is \emph{detected}, not hard-coded: the
JEDEC density byte is a power of two, so \texttt{Capacity\_Bytes} turns
\texttt{EF\,40\,19} into 32\,MB, and because each layer above sizes itself from the
one below, fitting a different W25Q part (8/16/32/64\,MB) needs no code change.
And the device record \emph{describes itself}: its host, its own bit clock, and its
chip select. The select is its own GPIO, not the SPI host's single hardware
\texttt{CS0}: the flash names its select pin and the SPI driver drives it
active-low, held across the whole command (so a streamed read can keep \texttt{CS}
asserted across many transfers). This is what lets the flash share the bus
with another device on \texttt{CS0}, each clocked its own way. Setting
\texttt{Clock\_Hz} and \texttt{CS\_Pin} is all it takes; a select that is not one
plain GPIO (a 3:8 decoder, an I/O-expander line) instead supplies a \texttt{CS\_CB}
callback. The host itself is brought up once with \texttt{ESP32S3.SPI.Setup} +
\texttt{Configure\_Pins} (the shared \texttt{Sclk}/\texttt{Mosi}/\texttt{Miso}).

\begin{lstlisting}
with ESP32S3.W25Q; use ESP32S3.W25Q;
--  Flash : ESP32S3.W25Q.Flash := (Host => ESP32S3.SPI.SPI2,
--    Clock_Hz => 8_000_000, CS_Pin => 21, others => <>);  --  8 MHz, CS on IO21
   ID      : JEDEC_ID;
   Mode_OK : Boolean;
begin
   Read_Identification (Flash, ID);          --  EF 40 19 on a W25Q256FV
   Initialize (Flash, Mode_OK);              --  enter 4-byte address mode
   if ID.Manufacturer = 16#EF# and then Capacity_Bytes (ID) /= 0 then
      null;                                  --  Capacity_Bytes (ID) = 32 MB
   end if;
\end{lstlisting}

\subsection{Block device and wear leveling}\index{W25Q\_Source}\index{block-erase}\index{ping-pong configuration}

NOR flash must be \emph{erased} (to all-ones, in 4\,KiB sectors) before it can be
re-programmed, so \texttt{W25Q\_Source} hides an erase-aware, write-through
512-byte sector interface over the chip; an optional block-erase capability lets
the layer above clear a whole 4\,KiB block in one operation.

\texttt{Block\_Dev.WL} is the wear-leveling FTL. It keeps an $O(1)$ state (a
single move counter) and maps a logical 4\,KiB block to a physical one by a
shifting offset, with one always-free ``hole''. Every few writes it advances the
hole by copying one block, so over time every logical block visits every physical
block and a hot block (an ext4 bitmap, say) never wears out one sector. The state
is written to two configuration blocks \emph{ping-pong} (alternating, each with a
higher sequence number and a CRC), so a power loss mid-update always leaves a
consistent earlier state; \texttt{Mount} picks the highest valid sequence.

That same ``highest sequence wins'' rule is why \texttt{Format} cannot just write a
fresh, low-sequence record and stop: a chip that was used before still holds an old
record with a far higher sequence in the \emph{other} ping-pong slot, and
\texttt{Mount} would then pick that stale slot and silently bring the volume back up
with its pre-format wear state, as though the format never happened.
\texttt{Format} therefore \emph{invalidates both} config slots first --- zeroing
them, so their magic and version no longer parse --- and only then writes the new
sequence-one record, so the fresh state is the only thing \texttt{Mount} can find.

The whole stack (erase-aware writes, the wear-level remap, the auto-detected size)
comes together in four calls:

\begin{lstlisting}
   BDW.Configure (Raw, Flash => Flash);      --  auto-size from the JEDEC id
   WL.Attach (Vol, BDW.Make (Raw'Access));   --  wear-leveling volume over it
   WL.Mount  (Vol, Formatted);               --  recover the saved wear state
   Dev := WL.Make (Vol'Access);              --  a 512-byte Block_Dev for ext4
\end{lstlisting}

\subsection{Formatting on-device: \texttt{Ext4.Mkfs}}\index{mkfs}\index{Ext4.Mkfs}\index{on-device format}

The external flash is not on the host flasher's path, and the filesystem reads and
writes but does not format, so to start from a blank chip the board makes the
filesystem itself. \texttt{ESP32S3.Ext4.Mkfs} is a minimal pure-Ada
\texttt{mkfs.ext4}: it writes a single-block-group ext4 (4\,KiB blocks, classic
block-mapped inodes, a root directory and \texttt{lost+found}), with an optional
4\,MiB JBD2 journal (\texttt{Journal => True}). It is the inverse of the read path,
validated against the host's \texttt{e2fsck}. From a blank chip to a mounted,
written filesystem is then:

\begin{lstlisting}
   --  ... build Dev as above, but WL.Format (Vol) for a fresh volume ...
   ESP32S3.Ext4.Mkfs.Format (Dev, Volume_Label => "FLASH", Journal => True);
   M.Open (Dev, Read_Only => False);
   N := M.Create_File ("/", "boot.txt");
   M.Write_File (N, (Character'Pos ('h'), Character'Pos ('i'), 10));
   M.Commit;                                 --  through the journal just created
\end{lstlisting}

\noindent The on-target examples \texttt{esp32s3\_w25q}, \texttt{esp32s3\_wl} and
\texttt{esp32s3\_ext4\_mkfs} run this stack end to end on the hardware.

\subsection{Writing in runs, not sectors}\index{Write\_Run}\index{erase amplification}

One more W25Q\_Source fact earns its space, because it is an 8$\times$
difference. The block layer's \texttt{Write\_Sector} is erase-aware but
per-sector: any change that is not purely clearing bits read-modify-writes the
whole 4\,KB erase block. Hand it the eight 512-byte sectors of one block, one
at a time, and it erases that block eight times --- eight times the wear, and at
$\sim$45\,ms an erase, eight times the wait. \texttt{Write\_Run} walks a run
of consecutive sectors one \emph{erase block} at a time instead: a block the
run covers completely is erased once and programmed; a block covered in part
takes one read-modify-write, not one per sector. Callers reach it through
\texttt{Block\_Dev.Write\_Sectors}, which uses the run primitive when the
device offers one and falls back to the per-sector loop when it does not.
Measured on a 32\,MB part behind a USB mass-storage function: a 64\,KiB
write fell from a per-sector 8\,kB/s crawl to wire speed.

\section{FAT16: the filesystem a PC can just mount}
\index{FAT16}\index{Fat16.Mkfs}\index{long filenames}

ext4 is the right home for the board's own data. It is exactly wrong for one
job: a volume a Windows PC must mount without drivers or questions.
\texttt{ESP32S3.Fat16} exists for that job --- a board exposes its flash over
USB mass storage, a PC drops files on it, the board reads them back. Interop
is the whole point, so the design follows what PCs expect rather than what
would be elegant.

The reader mounts an MBR-partitioned disk (the first FAT16 partition) or a
bare-boot-sector ``superfloppy'', walks long filenames as well as 8.3 names,
matches names case-insensitively because FAT's own rule does, and streams:
\texttt{Open} / \texttt{Read} / \texttt{Seek} over a caller-supplied buffer,
no heap. A \texttt{Read} that comes back short means end of file, never a
transient.

\texttt{Fat16.Mkfs.Format} lays down an empty volume so a blank chip appears
to the first PC as a usable drive rather than as ``you need to format this
disk''. Two facts drive its geometry. Windows expects a USB drive to carry a
partition table with the filesystem inside partition~1, so that is the
default. And the backing medium is NOR flash with a 4\,KB erase unit, so
clusters are 4\,KB and the data area starts on an erase-block boundary --- a
cluster never straddles two erase blocks, and a cluster write costs the block
layer one erase. Formatting is destructive and unconditional; deciding
\emph{whether} is the caller's job, and the honest test for ``blank'' is the
medium's own bytes (erased NOR reads all ones), not a missing signature ---
a volume the reader cannot parse looks exactly like no volume at all.

The host harness (\texttt{test/fat16\_host}) checks the pair three ways
against reference tooling: volumes written by \texttt{dosfstools}, volumes
written by Python, and volumes written by \texttt{Mkfs} itself, cross-read
and \texttt{fsck}-checked. One application built on this SDK runs the full
loop: it formats its own flash, serves it to a PC as a USB drive, and later
reads back the firmware images the PC left there.

